\section{What trial runs changed before the experiment}\label{sec:changes}

Before the real runs I made short trial runs to check that the set-up works. They led me to
change three things without asking you (points 1, 3 and 4) and showed one thing I had not
expected (point 2). The rule itself is untouched. But point 3 changes how a model is graded and
point 4 changes when a run may be called flat, and the rule depends on both.

\textbf{1. The speed of learning.} A setting called the learning rate decides how strongly the
numbers are changed at each step. My first value was too cautious: with a value five times
higher, width 8 learned the 16 animals about six times sooner. I chose the value using the
16-animal language only, so that I could not, even without meaning to, pick one that gave a
particular result on the large languages. Every width and every language used that one value.

\textbf{2. In the trial runs the smallest model learned no facts.} Width 4 learned the shapes
of the sentences, but its score was still 0 after 360 sets, even with only 16 animals (section
\ref{sec:four}). A width that fails condition (a) cannot count in favor of your idea, and
width 32 cannot either, because condition (c) needs a model twice as wide. So in the first
round the rule could only be met for widths 8 and 16.

\textbf{3. The score wobbled, so I grade a steadied copy of the model.} With the higher
learning rate, the numbers of a small model on a large language keep shifting a little at every
step, and the score shifts with them. In a trial run of width 8 on 1,024 animals, four gradings
in a row gave 97.6, 96.5, 91.9 and 100.0: a swing of 8 points, too much for a rule that turns on
10. The standard cure is to keep, beside the model being trained, a second copy in which each
number is the average of what that number has been over the last few hundred steps, and to
grade that copy. The training itself is not changed. To see what difference this makes I
repeated one such run without it. The model learned the language just the same, but its score
went on swinging: 96.1, 99.8 and 95.8 at three gradings in a row, so it did not pass the test
for ``solved'' (97 or more twice running) in the 397 sets I gave it. The steadied copy had
scored 100 and 100 at the first two of those gradings. One side effect of the copy: it is always a little
behind the model being trained, by roughly ten sets in the first round.

\textbf{4. A model can look finished before it has begun.} A model first learns the sentence
shapes, then stands still for a while with a score of zero, and only then starts to learn
facts. The ``flat'' test could not tell that pause from a model at its limit: it ended one trial
run, width 16 on 1,024 animals, as flat with a score of 0.4, although that width later learned
every fact of that language. A false ``flat'' on a large language is exactly what condition (b)
looks for, so this could have produced false support for your idea. Two conditions now
guard against it: a run cannot be called flat in its first 2,000 training steps (added when the
programs were reviewed), nor while its score is below 5 (added after this trial run).
A run that never starts to learn facts now ends as ``not finished''. The cost: a model that
truly cannot learn anything of a large language now looks the same as one that needs more time. Even so
repaired, the test was fooled once more, by a later and longer pause (section
\ref{sec:round32}).

\textbf{Checking the programs.} The programs that write the sentences and grade the models
were checked before use by other copies of Claude, the AI that wrote this report, each working
from the programs alone. They confirmed that every sentence is true, that each set holds
exactly the agreed number of each fact, and that a made-up ``model'' which knows everything is
graded 100 and one which knows nothing is graded 0.
